% ****** Start of file apssamp.tex ******
%
%   This file is part of the APS files in the REVTeX 4.2 distribution.
%   Version 4.2a of REVTeX, December 2014
%
%   Copyright (c) 2014 The American Physical Society.
%
%   See the REVTeX 4 README file for restrictions and more information.
%
% TeX'ing this file requires that you have AMS-LaTeX 2.0 installed
% as well as the rest of the prerequisites for REVTeX 4.2
%
% See the REVTeX 4 README file
% It also requires running BibTeX. The commands are as follows:
%
%  1)  latex apssamp.tex
%  2)  bibtex apssamp
%  3)  latex apssamp.tex
%  4)  latex apssamp.tex
%
\documentclass[%
 reprint,
%superscriptaddress,
groupedaddress,
%unsortedaddress,
% runinaddress,
%frontmatterverbose, 
% preprint,
%preprintnumbers,
%nofootinbib,
%nobibnotes,
%bibnotes,
 amsmath,amssymb,
 aps,
% pra,
%prb,
%rmp,
%prstab,
%prstper,
%floatfix,
]{revtex4-2}

\usepackage{graphicx}% Include figure files
\usepackage{dcolumn}% Align table columns on decimal point
\usepackage{bm}% bold math

\usepackage{enumitem,amssymb}
\newlist{checklist}{itemize}{2}
\setlist[checklist]{label=$\square$}
\usepackage{pifont}
\newcommand{\cmark}{\ding{51}}%
\newcommand{\xmark}{\ding{55}}%
\newcommand{\done}{\rlap{$\square$}{\raisebox{2pt}{\large\hspace{1pt}\cmark}}%
\hspace{-2.5pt}}
\newcommand{\wontfix}{\rlap{$\square$}{\large\hspace{1pt}\xmark}}

\usepackage{xcolor}
\newcommand{\DB}[1]{\textcolor{blue}{Dima: #1}}
\newcommand{\YZ}[1]{\textcolor{purple}{Yuzhe: #1}}
\newcommand{\Copilot}[1]{\textcolor{brown}{Copilot: #1}}

\usepackage{ulem}

\begin{document}

\preprint{APS/123-QED}

\title{Five Rules of Good Scientific Communication}

\author{Yuzhe Zhang}
\email{yuhzhang@uni-mainz.de}
\affiliation{Johannes Gutenberg-Universit{\"a}t Mainz, 55128 Mainz, Germany}
\affiliation{Helmholtz Institute Mainz, 55099 Mainz, Germany}
\affiliation{GSI Helmholtzzentrum für Schwerionenforschung GmbH, 64291 Darmstadt, Germany}


\author{Dmitry Budker}
\email{budker@uni-mainz.de}
\affiliation{Johannes Gutenberg-Universit{\"a}t Mainz, 55128 Mainz, Germany}
\affiliation{Helmholtz Institute Mainz, 55099 Mainz, Germany}
\affiliation{GSI Helmholtzzentrum für Schwerionenforschung GmbH, 64291 Darmstadt, Germany}
\affiliation{Department of Physics, University of California, Berkeley, CA 94720-7300, United States of America}


\date{September 9, 2026}

\begin{abstract}
Some common rules for transparent scientific communications in all forms are presented. Following these rules can help avoid the basic but also frequent mistakes.  

\end{abstract}


\maketitle

\section{Introduction}

Excellent scientific communication demands excellent understanding of science in the first place, and then demands good communication skills. 
% Scientific communications take forms of articles, colloquium, keynote speech, elevator chat, and etc. 
Previous education notes of the authors and the collaborators have covered communication media including article\,\cite{Budker2008writing, Budker_Kimball_2016_collaborative}, poster\,\cite{Budker_Trabesinger_2017_Poster}, and viewgraph\,\cite{Zhang_Budker_2024_viewgraphs}. 
In this note, we do not focus on a specific form or medium of communications, but highlight some basic rules for communication in general. 
Each of the rules corresponds to a basic but unfortunately common mistake. 

\section{The rules}

\begin{itemize}

\item \textbf{One thing at a time.} Do not challenge the concentration of the audience. One common mistake is introducing two topics within one viewgraph of a presentation. 

\item \textbf{Good answers are direct.} Who $\rightarrow$ someone; what $\rightarrow$ something; where $\rightarrow$ someplace; when $\rightarrow$ sometime; why $\rightarrow$ some reason; how $\rightarrow$ a description; Yes/no question $\rightarrow$ yes/no. 
% Start your answer by a short and direct response. 

\item \textbf{Introduce every new element;} Here ``element'' includes theorems, abbreviations, adjectives, and often yourself. 
The audience's knowledge of the elements is always overestimated. 
Repeating definitions of new elements is beneficial to 
\begin{itemize}
    \item reinforce audience's memory;
    % \item make sure the audience is synchronized with you;
    \item slow down the flow of information, which is usually too fast and overwhelming; 
    \item leave some room for snoozing.
\end{itemize}

% \item \textbf{People are only smart with numbers smaller than 20.} Do not distract your audience by making them do math in heart. 

\item \textbf{Example helps.} An example is worth a thousand sentences of explanation. 


\item \textbf{Highlights are made by contrast.} For example, relative magnitudes are favored over absolute ones in demonstrating improvements. 

% \item \textbf{Left to right / low to high corresponds naturally to small to large.} Better not to invert this relationship.

% \item \textbf{Humans will naturally compare plots with the same $x$- and $y$-axis labels.}

\end{itemize}


\section{Concluding remarks}

The rules presented can guarantee transparent communication. 
Following these rules are not as easy as it seems. 
Artificial intelligence (AI) can help you check how well your article or speech follows the rules. 


% \begin{acknowledgments}

% \end{acknowledgments}

\bibliography{references.bib}% Produces the bibliography via BibTeX.

\end{document}
%
% ****** End of file apssamp.tex ******
